% Part: methods
% Chapter: proofs
% Section: example-1

\documentclass[../../../include/open-logic-section]{subfiles}

\begin{document}

\olfileid{mth}{prf}{ex1}

\olsection{یک مثال}

نخستین مثال ما واقعیت سادهٔ زیر دربارهٔ اجتماع و اشتراک مجموعه‌هاست.
این مثال بازکردن تعریف‌ها، اثبات عطف‌ها و گزاره‌های کلی، و اثبات به
تفکیک حالت‌ها را نشان می‌دهد.

\begin{prop}
برای هر سه مجموعهٔ $A$، $B$ و $C$،
$A \cup (B \cap C) = (A \cup B) \cap (A \cup C)$.
\end{prop}

بیایید اثباتش کنیم!{}

\begin{proof}
می‌خواهیم نشان دهیم که برای هر سه مجموعهٔ $A$، $B$ و $C$،
$A \cup (B \cap C) = (A \cup B) \cap (A \cup C)$.
\begin{quote}
ابتدا تعریف «$=$» را در صورت حکم باز می‌کنیم. به یاد آورید که برای
اثبات یکسان‌بودن دو مجموعه باید نشان دهیم آن‌ها !!{element}های یکسانی
دارند. یعنی همهٔ !!{element}های $A \cup (B \cap C)$، !!{element}های
$(A \cup B) \cap (A \cup C)$ نیز هستند، و برعکس. «برعکس» یعنی هر
!!{element} از $(A \cup B) \cap (A \cup C)$ نیز باید !!a{element} از
$A \cup (B \cap C)$ باشد. پس با بازکردن تعریف می‌بینیم که باید یک عطف
را اثبات کنیم. این را ثبت می‌کنیم:
\end{quote}
بنا بر تعریف، $A \cup (B \cap C) = (A \cup B) \cap (A \cup C)$ اگر و
تنها اگر هر !!{element} از $A \cup (B \cap C)$، !!a{element} از
$(A \cup B) \cap (A \cup C)$ نیز باشد، و هر !!{element} از
$(A \cup B) \cap (A \cup C)$، !!a{element} از $A \cup (B \cap C)$ باشد.
\begin{quote}
چون این گزاره یک عطف است، باید هر مؤلفهٔ عطف را جداگانه اثبات کنیم.
از مؤلفهٔ نخست آغاز کنیم: نشان دهیم هر !!{element} از
$A \cup (B \cap C)$، !!a{element} از $(A \cup B) \cap (A \cup C)$ نیز
هست.

این یک گزارهٔ کلی است؛ پس یک !!{element} دلخواه از $A \cup (B \cap C)$
را در نظر می‌گیریم و نشان می‌دهیم که باید !!a{element} از
$(A \cup B) \cap (A \cup C)$ نیز باشد. متغیری برای نامیدن این
!!{element} دلخواه برمی‌گزینیم، مثلاً~$z$. اثبات چنین ادامه می‌یابد:
\end{quote}
نخست اثبات می‌کنیم که هر !!{element} از $A \cup (B \cap C)$،
!!a{element} از $(A \cup B) \cap (A \cup C)$ نیز هست. فرض کنید
$z \in A \cup (B \cap C)$. باید نشان دهیم
$z \in (A \cup B) \cap (A \cup C)$.
\begin{quote}
اکنون وقت آن است که تعریف‌های $\cup$ و~$\cap$ را باز کنیم. برای نمونه،
تعریف $\cup$ چنین است:
$A \cup B = \Setabs{z}{z \in A \text{ یا } z \in B}$.
وقتی تعریف را بر «$A \cup (B \cap C)$» اعمال می‌کنیم، نقش «$B$» در
تعریف را اکنون «$B \cap C$» بازی می‌کند؛ بنابراین
$A \cup (B \cap C) = \Setabs{z}{z \in A \text{ یا } z \in B \cap C}$.
پس فرض $z \in A \cup (B \cap C)$ به این معناست که
$z \in \Setabs{z}{z \in A \text{ یا } z \in B \cap C}$. همچنین
$z \in \Setabs{z}{\dots z\dots}$ اگر و تنها اگر \dots $z$ \dots؛ یعنی
در این مورد، $z \in A$ یا $z \in B \cap C$.
\end{quote}
بنا بر تعریف $\cup$، یا $z \in A$ یا $z \in B \cap C$.
\begin{quote}
چون این یک فصل است، اثبات به تفکیک حالت‌ها سودمند خواهد بود. دو حالت
را در نظر می‌گیریم و نشان می‌دهیم که در هر یک نتیجهٔ مورد نظر، یعنی
«$z \in (A \cup B) \cap (A \cup C)$»، برقرار است.
\end{quote}
حالت ۱: فرض کنید $z \in A$.
\begin{quote}
از فرض‌هایمان چیز زیادی برای ادامه در دست نیست. پس به نتیجه‌ای نگاه
کنیم که باید اثبات شود. می‌خواهیم نشان دهیم
$z \in (A \cup B) \cap (A \cup C)$. بر پایهٔ تعریف $\cap$، برای نشان
دادن این امر باید نشان دهیم $z$ هم در $(A \cup B)$ و هم در
$(A \cup C)$ است. اما $z \in A \cup B$ اگر و تنها اگر $z \in A$ یا
$z \in B$، و از پیش، به‌عنوان فرض حالت~۱، $z \in A$ را داریم. با همان
استدلال و با جایگزین‌کردن $C$ به جای $B$، $z \in A \cup C$. این
استدلال را در جهت معکوس پیمودیم؛ اکنون آن را در جهت لازم برای اثبات
ثبت می‌کنیم.
\end{quote}
چون $z \in A$، داریم $z \in A$ یا $z \in B$؛ پس بنا بر تعریف~$\cup$،
$z \in A \cup B$. به همین ترتیب، $z \in A \cup C$. اما بنا بر تعریف
~$\cap$، این یعنی $z \in (A \cup B) \cap (A \cup C)$.
\begin{quote}
این کار حالت نخست اثبات به تفکیک حالت‌ها را کامل می‌کند. اکنون باید
در حالت دوم، که $z \in B \cap C$ است، نتیجه را به دست آوریم.
\end{quote}
حالت ۲: فرض کنید $z \in B \cap C$.
\begin{quote}
باز هم با اشتراک دو مجموعه سروکار داریم. تعریف~$\cap$ را به کار می‌بریم:
\end{quote}
چون $z \in B \cap C$، بنا بر تعریف~$\cap$، $z$ باید !!a{element} از هر
دو مجموعهٔ $B$ و $C$ باشد.
\begin{quote}
وقت آن است که دوباره به نتیجه بنگریم. باید نشان دهیم $z$ هم در
$(A \cup B)$ و هم در $(A \cup C)$ است. پاسخ باز هم بی‌درنگ به دست می‌آید.
\end{quote}
چون $z \in B$، داریم $z \in (A \cup B)$. چون $z \in C$، همچنین
$z \in (A \cup C)$. پس $z \in (A \cup B) \cap (A \cup C)$.
\begin{quote}
اینجا دوباره تعریف‌های $\cup$ و $\cap$ را به کار بردیم؛ اما چون پیش‌تر
این تعریف‌ها را یادآوری کرده‌ایم و نشان داده‌ایم که عضویت $z$ در یکی
از دو مجموعه، عضویت آن در اجتماعشان را نتیجه می‌دهد، لازم نیست کارمان
را با همان تفصیل بیان کنیم.

حالت دوم اثبات به تفکیک حالت‌ها کامل شد؛ بنابراین اکنون می‌توانیم
نخستین نتیجه را اعلام کنیم.
\end{quote}
پس اگر $z \in A \cup (B \cap C)$، آنگاه
$z \in (A \cup B) \cap (A \cup C)$.
\begin{quote}
اکنون فقط باید جهت دیگر را نشان دهیم: هر !!{element} از
$(A \cup B) \cap (A \cup C)$، !!a{element} از $A \cup (B \cap C)$ است.
مانند پیش، این گزارهٔ کلی را با فرض‌کردن یک !!{element} دلخواه از
مجموعهٔ نخست و نشان‌دادن عضویت آن در مجموعهٔ دوم اثبات می‌کنیم. آنچه
می‌خواهیم انجام دهیم چنین است.
\end{quote}
اکنون فرض کنید $z \in (A \cup B) \cap (A \cup C)$. می‌خواهیم نشان دهیم
$z \in A \cup (B \cap C)$.
\begin{quote}
اکنون از فرض $z \in (A \cup B) \cap (A \cup C)$ کار را پیش می‌بریم.
امیدواریم استفاده از همان~$z$ که در بخش نخست اثبات به کار رفت گیج‌کننده
نباشد. با پایان آن بخش، همهٔ فرض‌هایی که در آن ساخته بودیم دیگر نافذ
نیستند؛ پس اکنون می‌توانیم فرض‌های تازه‌ای دربارهٔ $z$ بسازیم. اگر این
موضوع گیج‌کننده است، در ادامه $z$ را با متغیری دیگر جایگزین کنید.

بنا بر تعریف~$\cap$ می‌دانیم $z$ هم در $A \cup B$ و هم در $A \cup C$
است. با تعریف $\cup$ می‌توانیم این را بیشتر باز کنیم: یا $z \in A$ یا
$z \in B$، و نیز یا $z \in A$ یا $z \in C$. این دوباره شبیه اثبات به
تفکیک حالت‌هاست، جز آنکه «و» موضوع را پیچیده می‌کند. شاید تصور کنید
این سخن سه امکان می‌دهد: $z$ یا در $A$ است یا در $B$ یا در $C$. اما
این اشتباه است. باید دقت کنیم؛ پس هر فصل را به نوبت بررسی می‌کنیم.
\end{quote}
بنا بر تعریف $\cap$، $z \in A \cup B$ و $z \in A \cup C$. بنا بر
تعریف $\cup$، یا $z \in A$ یا $z \in B$. حالت‌ها را جدا می‌کنیم.
\begin{quote}
چون اکنون بر فصل نخست تمرکز داریم، هنوز فصل دوم را که از بازکردن
$A \cup C$ به دست می‌آید ننوشته‌ایم. در واقع هنوز به آن نیاز نداریم.
حالت نخست $z \in A$ است، و هر !!a{element} از یک مجموعه، !!a{element}
از اجتماع آن مجموعه با هر مجموعهٔ دیگر نیز هست. پس حالت~۱ آسان است:
\end{quote}
حالت ۱: فرض کنید $z \in A$. در نتیجه $z \in A \cup (B \cap C)$.
\begin{quote}
اکنون به حالت دوم، یعنی $z \in B$، می‌پردازیم. اینجا $\cup$ دوم را
باز می‌کنیم و بار دیگر اثبات به تفکیک حالت‌ها انجام می‌دهیم:
\end{quote}
حالت ۲: فرض کنید $z \in B$. چون $z \in A \cup C$، یا $z \in A$ یا
$z \in C$. باز هم حالت‌ها را جدا می‌کنیم:

حالت ۲الف: $z \in A$. پس باز هم $z \in A \cup (B \cap C)$.
\begin{quote}
این مورد کمی نامعمول بود. در واقع در این حالت به فرض $z \in B$ نیاز
نداشتیم، اما اشکالی ندارد.
\end{quote}
حالت ۲ب: $z \in C$. آنگاه $z \in B$ و $z \in C$؛ پس
$z \in B \cap C$ و در نتیجه $z \in A \cup (B \cap C)$.
\begin{quote}
بدین ترتیب هر دو اثبات به تفکیک حالت‌ها پایان می‌یابد و نیمهٔ دوم نیز
کامل می‌شود.
\end{quote}
پس اگر $z \in (A \cup B) \cap (A \cup C)$، آنگاه
$z \in A \cup (B \cap C)$.
\end{proof}

\end{document}
